% ==============================================================================
% CAPÍTULO 3 - REFERÊNCIAL TEÓRICO
% ==============================================================================
    
\section{REFERÊNCIAL TEÓRICO (opcional)} 
\label{sec:referencial-teorico}

Nesta seção, deve-se apresentar o estado da arte da pesquisa, discutindo e 
comparando trabalhos recentes que possuam objetivos, problemas ou metodologias 
semelhantes ao tema abordado. O texto deve ir além da simples descrição, 
estabelecendo um diálogo entre as obras citadas para destacar as lacunas 
que o seu trabalho pretende preencher.

Escreva um parágrafo introdutório apresentando os critérios de busca ou os 
grandes eixos temáticos que agrupam as pesquisas selecionadas. Utilize subseções 
para organizar os trabalhos por similaridade (ex: Abordagens Baseadas em IA, 
Estudos de Caso em Engenharia de Software). É fundamental que, ao final de cada 
análise, fique claro como a sua proposta se diferencia ou avança em relação aos
artigos apresentados.

\textbf{Orientações para Citação}

Abaixo, estão disponibilizados exemplos de como realizar citações no texto
utilizando comandos do LaTeX. Lembre-se de que todas as fontes devem estar
cadastradas no seu arquivo de referências (\texttt{.bib}) para que a lista
final do trabalho seja gerada automaticamente.

\textbf{Exemplos de citações:}

Use o \texttt{\textbackslash citeonline\{\}} quando você quiser citar o nome
do autor diretamente no seu texto (o LaTeX deixará apenas o ano entre parênteses).

\textbf{Exemplo:} Segundo \citeonline{almeida2022}, a escolha correta de algoritmos
de otimização é fundamental para melhorar o desempenho e a velocidade de
resposta em sistemas distribuídos.

Use o \texttt{\textbackslash cite\{\}} quando fizer uma afirmação e quiser
indicar a fonte apenas ao final (o LaTeX colocará o nome e o ano entre
parênteses).

\textbf{Exemplo:} A utilização de sistemas distribuídos exige uma análise
detalhada sobre como os algoritmos de otimização podem reduzir o tempo de
processamento \cite{almeida2022}.

\textbf{Citações usadas para gerar as referências de exemplo:}\\
Estes exemplos são utilizados para mostrar como diferentes tipos de
fontes são apresentados ao final do documento nas referências.

Verifique os campos necessários no arquivo ``referencias.bib'' para que 
cada referência seja apresentada corretamente na seção de Referências: 

Dissertação, \cite{almeida2022}\\
Trabalho publicado em Conferência, \cite{cbsoft2023}\\
Página web, \cite{ferreira2020}\\
Manual, \cite{latex2023}\\
Capitulo de livro, \cite{mendes2019}\\
Relatório técnico, \cite{oliveira2020}\\
Livro, \cite{linden2012}\\
Tese de doutorado, \cite{rodrigues2021}\\
Artigo científico, \cite{silva2023}\\
Livro de Conferência, \cite{souza2022}